\documentclass[11pt]{article}
\usepackage[margin=1in]{geometry}
\usepackage[T1]{fontenc}
\usepackage[utf8]{inputenc}
\usepackage{amsmath,amssymb,amsthm}
\usepackage{hyperref}
\usepackage{listings}
\setlength{\emergencystretch}{4em}
\hypersetup{hidelinks}
\lstset{basicstyle=\ttfamily\small,breaklines=true,columns=fullflexible,keepspaces=true}
\newtheorem{theorem}{Theorem}
\newtheorem{lemma}[theorem]{Lemma}
\title{A square-root counterexample to the proposed Kingman rate\\Conjecture 00000008482}
\author{AlyciaBHZ, on behalf of the Omega Institute}
\date{5 October 2026}
\begin{document}
\maketitle
\begin{abstract}
The deterministic stationary subadditive family $X_{m,n}=\sqrt{n-m}$ has
unit one-step costs, growth at most $n-m$, and successive increments at most
one. Its limiting normalized expectation is zero, and its convergence error
is exactly $n^{-1/2}$. This is not $O(\log n/n)$, refuting the proposed
universal upper bound under the reading of bounded costs or bounded
increments. We explicitly distinguish the alternative reading of a single
bound over all interval lengths: under that reading every error is
$O(1/n)$, which instead refutes the asserted logarithmic optimality.
Both readings are checked in Lean 4 with Mathlib: the square-root
counterexample refutes the upper bound, while uniform boundedness forces
zero limit and error $o(\log n/n)$, refuting logarithmic sharpness.
\end{abstract}

\section{Statement and interpretation}
Problem \textbf{00000008482} states, verbatim:
\begin{quote}
Definition: The convergence error of Kingman's subadditive ergodic theorem: the decay order of |n\textasciicircum\{-1\}E[X\_n] - gamma| in n. Conjecture: The universal upper bound for bounded subadditive families is O(log n/n) (a logarithmic harmonic law); the bound cannot be improved, extremal examples being realized by discretization deviations of Sturmian balanced words (a Sturmian extremal law); and in the unbounded case the error can decay arbitrarily slowly (an arbitrary slowness law). (Kingman Sturmian extremal logarithmic decay)
\end{quote}
We identify $X_n$ with $X_{0,n}$ of a stationary two-index subadditive
process. The conjecture does not define ``bounded.'' Our primary reading
is bounded local costs or bounded increments, as used for subadditive
processes that can accumulate costs over an interval. We require the
following explicit, simultaneous restrictions:
\[
0\le X_{m,n}\le n-m,\qquad X_{m,m+1}=1,\qquad
|X_{m,n+1}-X_{m,n}|\le1\quad(m\le n).
\]
The first condition bounds all normalized interval costs, and also implies
that each $X_{m,n}$ is a bounded random variable. The second uniformly
bounds one-step costs; the third uniformly bounds successive increments.
These are substantial restrictions, imposed together on the counterexample.

There is another literal reading: one constant $B$ bounds $|X_{m,n}|$ for
every interval, independently of its length. Our example does \emph{not}
satisfy that condition, since $X_{0,n}=\sqrt n$ is unbounded as $n$ varies.
We do not claim otherwise. Section~\ref{sec:global} addresses that reading
directly. Under it the first conjunct is true, but the claimed optimality
of the logarithmic rate is false. Thus we neither silently redefine
boundedness nor claim that this particular counterexample refutes the
first conjunct under every possible meaning of ``bounded.''

\section{The deterministic process}
For integers $0\le m\le n$, set
\[X_{m,n}=\sqrt{n-m}.\]
On a probability space this means the constant random variable with that
value. In particular all expectations exist and
$a_n=\mathbb E[X_{0,n}]=\sqrt n$. One can take a one-point probability
space with the identity shift; this shift is measure preserving and
ergodic. Thus the example belongs even to the ergodic setting, not merely
the stationary setting. Joint stationarity holds pathwise, because
\[X_{m+k,n+k}=\sqrt{(n+k)-(m+k)}=X_{m,n}.\]

\begin{lemma}
For $a,b\ge0$, $\sqrt{a+b}\le\sqrt a+\sqrt b$.
\end{lemma}
\begin{proof}
Both sides are nonnegative and
$(\sqrt a+\sqrt b)^2=a+b+2\sqrt{ab}\ge a+b$.
Squaring therefore proves the inequality.
\end{proof}
For $l\le m\le n$, put $a=m-l$, $b=n-m$. The lemma yields
\[X_{l,n}\le X_{l,m}+X_{m,n}.\]
Nonnegativity is immediate. For an integer $d=n-m$, either $d=0$ or
$d\ge1$, and in both cases $\sqrt d\le d$. This proves the linear bound.
The one-step cost is $\sqrt1=1$. Monotonicity and the lemma with $b=1$
give
\[0\le\sqrt{d+1}-\sqrt d\le1.\]
Hence the stated absolute increment bound holds for every interval.
All negative-part integrability assumptions of Kingman's theorem are
automatic, and $\mathbb E[X_{0,1}]=1<\infty$.

\section{Limit and failure of the asserted rate}
For every $n\ge1$,
\[\frac{a_n}{n}=\frac{\sqrt n}{n}=\frac1{\sqrt n}.
\]
Consequently
\[\gamma=\lim_{n\to\infty}\frac{a_n}{n}
  =\inf_{n\ge1}\frac{a_n}{n}=0.
\]
The equality with the infimum follows because all terms are nonnegative
and converge to zero. Thus the convergence error is exactly
\[\left|\frac{a_n}{n}-\gamma\right|=\frac1{\sqrt n}.
\]

\begin{theorem}\label{thm:main}
This error is not $O(\log n/n)$. There is no per-family big-O bound,
and therefore no bound with a constant uniform over families.
\end{theorem}
\begin{proof}
For $n\ge2$, the ratio of the error to the proposed scale is
\[\frac{1/\sqrt n}{\log n/n}=\frac{\sqrt n}{\log n}\longrightarrow\infty.
\]
For completeness the elementary logarithm estimate is enough to establish
this divergence: for $x\ge1$,
\[\log x=\int_1^x\frac{dt}{t}
   \le\int_1^x t^{-3/4}\,dt=4(x^{1/4}-1)\le4x^{1/4}.
\]
It follows that $\log x/\sqrt x\to0$ and the displayed reciprocal diverges.
Equivalently, a purported eventual inequality
$\sqrt n\le C\log n$ is incompatible with $\log n=o(\sqrt n)$.
The latter is the form used in Lean. Multiplying the hypothetical big-O
bound by $n$ gives $\sqrt n=O(\log n)$; composing with
$\log n=o(\sqrt n)$ gives $\sqrt n=o(\sqrt n)$, impossible since
$\sqrt n>0$ eventually. This argument allows both the constant and the
starting index to depend on the family, so it disproves the weaker
per-family reading as well as the uniform reading.
\end{proof}
The original conjecture is a conjunction. Under the bounded-costs reading,
Theorem~\ref{thm:main} refutes its first conjunct, so no separate
disproof of the Sturmian assertion or the arbitrary-slowness assertion
is required. We make no claim here that the arbitrary-slowness assertion
is false.

\section{Robustness to the reading of 'bounded'}\label{sec:global}
Suppose instead that a single $B<\infty$ satisfies
$|X_{m,n}(\omega)|\le B$ for all intervals almost surely.
Then $|\mathbb E[X_{0,n}]|\le B$, so $a_n/n\to0$, and
\[\gamma=0,\qquad |a_n/n-\gamma|\le B/n.
\]
The same reasoning applies if the single bound is imposed directly on
the expectations. Since $1/n=o(\log n/n)$, the logarithmic order is
strictly improvable throughout this class. In particular no example in
this class, Sturmian or otherwise, can attain a non-improvable
$\log n/n$ order.

The $1/n$ order itself can be attained. On a one-point probability space
take $Y_{m,n}=1$ when $m<n$ and $Y_{m,m}=0$. It is stationary and globally
bounded by one. For $l\le m\le n$, if $l=n$ the subadditivity inequality
is zero on both sides. If $l<n$, at least one of $l<m$ or $m<n$ holds,
so the right side is at least one, proving subadditivity. Now
$\mathbb E[Y_{0,n}]=1$ for $n\ge1$ and the error is exactly $1/n$.
The uniform-boundedness argument is now also formalized. In Lean,
\texttt{UniformlyBounded X} means exactly
\[
\exists C\in\mathbb R,\quad \forall m,n\in\mathbb N,\quad |X_{m,n}|\le C.
\]
The theorem \texttt{uniformlyBounded\_limit} proves both convergence of
$X_{0,n}/n$ to zero and \texttt{gamma X = 0};
\texttt{uniformlyBounded\_error\_isLittleO} proves
\[
\left|X_{0,n}/n-\gamma X\right|=o(\log n/n).
\]
These results require only uniform boundedness. The little-o proof uses
$X_{0,n}=O(1)$ and $1=o(\log n)$, then multiplies both comparison functions
by $1/n$ and takes the absolute value. Thus the first upper-bound clause
is true under this reading, but its claimed sharpness is false.

For this reading the admissibility predicate
\texttt{UniformlyBoundedStationarySubadditive} contains exactly translation
stationarity, subadditivity on ordered triples, and
\texttt{UniformlyBounded}. It imposes no nonnegativity, normalization,
unit one-step cost, or separate increment restriction. The sharpness
clause \texttt{SharpnessClauseUniform} asserts the existence of such a
family whose error is \emph{not} $o(\log n/n)$.
The checked theorem \texttt{sharpness\_false\_uniform} rules out every such
witness. The combined theorem is
\begin{lstlisting}
theorem conjecture_00000008482_false_both_readings :
    Not UniversalUpperBound /\ Not SharpnessClauseUniform
\end{lstlisting}
Under the linear-growth or bounded-increment reading the first conjunct
fails, witnessed by the square-root family. Under the uniformly bounded
reading the second conjunct fails for every admissible family. Either way
the original conjunction is false. The $1/n$ witness just described is a
written example; the universal limit and little-o conclusions, sharpness
refutation, and combined result are all Lean theorems.

\section{Formal statement and semantic audit}
The project pins Lean \texttt{leanprover/lean4:v4.33.0} and Mathlib
\texttt{v4.33.0}. The Mathlib revision is:
\begin{lstlisting}
db584cd6d46c92f209a44c0f1c829460d327499d
\end{lstlisting}
The single source file is \texttt{lean4/KingmanCounterexample.lean}.
All declarations are in namespace \texttt{KingmanCounterexample}.
\texttt{Family} abbreviates $\mathbb N\to\mathbb N\to\mathbb R$.
The value of \texttt{sqrtFamily m n} is
\texttt{Real.sqrt (n - m : Nat)}. Only ordered intervals are relevant;
natural subtraction simply supplies harmless values elsewhere.

The predicate \texttt{BoundedStationarySubadditive X} has exactly five
fields: subadditivity on ordered triples, translation stationarity,
the nonnegative linear bound, unit one-step costs, and the absolute
successive-increment bound. Requiring exact unit costs narrows the class
further, making the negative result stronger: the broader conjectured
universal bound would imply this restricted universal bound.

\texttt{gamma X} is defined as
\begin{lstlisting}
limUnder atTop (fun n : Nat => X 0 n / (n : Real))
\end{lstlisting}
This is Mathlib's limit operator. We do not presume convergence for an
arbitrary function: the universal clause explicitly requires convergence
to this value. Our witness satisfies this extra premise, as proved by
\texttt{sqrtFamily\_limit} and \texttt{sqrtFamily\_gamma}.
The universal clause is the following actual Lean definition (ASCII
spellings of mathematical symbols are used here):
\begin{lstlisting}
def UniversalUpperBound : Prop :=
  forall X : Family, BoundedStationarySubadditive X ->
    Tendsto (fun n : Nat => X 0 n / (n : Real))
      atTop (nhds (gamma X)) ->
    Asymptotics.IsBigO atTop
      (fun n : Nat => abs (X 0 n / (n : Real) - gamma X))
      (fun n => Real.log (n : Real) / (n : Real))

theorem conjecture_00000008482_false :
    Not UniversalUpperBound
\end{lstlisting}
The displayed definition is a typographical transcription of the
source's Unicode notation, with the same quantifiers and arguments.
Mathlib's \texttt{IsBigO} has its usual meaning: an eventual norm bound
with some real constant. The filter is the natural-number \texttt{atTop}.
Therefore this is a genuine asymptotic assertion on the full sequence,
not a finite sample or selected subsequence.

The supporting checked theorems establish the following:
\begin{itemize}
\item \texttt{sqrtFamily\_hypotheses}: all five family conditions.
\item \texttt{sqrtFamily\_error}: the error is the reciprocal square root.
\item \texttt{sqrtFamily\_limit} and \texttt{sqrtFamily\_gamma}: the limit exists
and equals zero.
\item \texttt{sqrtFamily\_infimum}: the infimum over lengths $n+1$ is zero,
covering every positive length without admitting length zero.
\item \texttt{sqrtFamily\_not\_bigO}: the error fails the asserted big-O bound.
\item \texttt{conjecture\_00000008482\_false}: the restricted universal clause
is false, using that explicit witness.
\item \texttt{uniformlyBounded\_limit}: uniform boundedness forces
convergence and \texttt{gamma X = 0}.
\item \texttt{uniformlyBounded\_error\_isLittleO}: every such error is
$o(\log n/n)$.
\item \texttt{sharpness\_false\_uniform}: the uniformly bounded sharpness
clause is false with exactly the three relevant admissibility conditions.
\item \texttt{conjecture\_00000008482\_false\_both\_readings}: the two
refutations hold together.
\end{itemize}
At length zero Lean uses totalized division by zero, giving zero; this
does not affect any \texttt{atTop} assertion or the positive-length infimum.
The deterministic formulation needs no expectation machinery: constant
random variables on any probability space have exactly these expectations.
The probabilistic embedding and one-point ergodicity explanation above
are elementary written arguments.

\section{Verification and attribution}
The Lean compilation succeeded with Lean and Mathlib \texttt{v4.33.0};
\texttt{verification.txt} records the compiler output and axiom audit. The audited main and supporting theorems use
only \texttt{propext}, \texttt{Classical.choice}, and \texttt{Quot.sound}.
There are no added axioms or incomplete proofs. The asymptotic step uses
Mathlib's proved theorem \texttt{isLittleO\_log\_rpow\_atTop} with
exponent $1/2$, composed with natural-number inclusion, and a proved
big-O/little-o incompatibility.

\texttt{verify.py} is an independent Python standard-library sanity check
of subadditivity on finite ordered triples, stationarity under shifts,
the bounds, the exact error formula numerically, and growing ratios on
geometrically separated indices. Floating-point checks are explicitly
supplementary; none is used as a proof of an infinite assertion. To
reproduce in a fresh environment:
\begin{lstlisting}
cd lean4
lake exe cache get && lake build
cd ..
python3 verify.py
pdflatex report.tex
pdflatex report.tex
\end{lstlisting}
No external project code is vendored. The source conjecture is available
in the competition repository at
\href{https://github.com/The-Last-Math-Competition/The-Last-Math-Competition/blob/main/conjectures/00000008482.md}{conjectures/00000008482.md}.

\medskip\noindent
Submitted by AlyciaBHZ on behalf of the Omega Institute (trureturing project,
\url{https://github.com/the-omega-institute/trureturing}).
\end{document}
